总长 \(L\) 的不可伸长柔链越过半径可忽略的无摩擦光滑定滑轮，两段竖直下垂，均未触底。用沿链的材料坐标 \(s\in[0,L]\) 标记左端起算的质元，线密度 \(\lambda(s)=\lambda_0(1+\alpha s/L)\)，其中 \(\alpha>-1\)。设左段长度 \(y\)，链运动时材料坐标 \(s=y\) 恰在滑轮顶点，右段长度 \(L-y\)。两段速度大小同为 \(|\dot y|\)，取左段增长为正，忽略弯曲和滑轮质量。
(1) 求总质量、左右两侧的瞬时重量之差，并正确写出 \(y\) 的运动方程。
(2) 以从静止的 \(y_0\) 出发为初值，给出 \(v^2(y)\) 的精确多项式表达，并指出允许的运动区间。
(3) 将抵达任意 \(y_1\) 的时间写成可数值计算的定积分；说明为何不能把不均匀链简单当作两团质量不变的质点。
(4) 确定唯一静止平衡位置及其线性稳定性，求微小偏移的指数增长率，并检验均匀链极限。
